\subsection{The parameter on the genus-two hyperelliptic quotient}

\begin{result}{B.4.2}{Proved}{The genus-two quotient has the same GT parameter}
Let $\alpha_F$ be the full genus-two outer action supplied by A.6 for
$F\in\widehat{\mathrm{GT}}$. Let $\iota$ be the hyperelliptic involution
and use the actual quotient
\[
 q:\widehat\Gamma_2\twoheadrightarrow
       G:=\widehat\Gamma_{0,[6]},\qquad
 \ker q=\langle\iota\rangle\simeq C_2.                 \tag{B.P1}
\]
The action induced by $\alpha_F$ on $G$ is the ordinary genus-zero
outer action with parameter $F$. In particular, for a chosen
representative $\gamma_F\in\operatorname{Aut}(G)$ of that ordinary
action, there is an actual $v\in\widehat\Gamma_2$ such that
\[
 q\circ\operatorname{Inn}(v^{-1})\circ\alpha_F
                    =\gamma_F\circ q.                 \tag{B.P2}
\]
Here a representative of $\alpha_F$ is understood in (B.P2).
This assertion uses the ordinary sewing action of A.6 and the
genus-zero classification; it invokes no universal higher-genus
quotient assertion.
\end{result}

\begin{proof}
Choose the positive five-chain twists $T_1,\ldots,T_5$ on a
genus-two surface. Their hyperelliptic quotients are the positive
adjacent branch half-twists $\sigma_1,\ldots,\sigma_5$, with
$q(T_i)=\sigma_i$. The Birman--Hilden presentation and its exact
central quotient are recorded in
\cite[Theorems 3.2--3.3 and the following central-extension display,
pp.~6--7]{CA_EndoKamada2016}. In genus two the five-chain twists
generate the whole mapping-class group. Completion of the discrete
central quotient is exact: every finite quotient of the quotient
group is precisely a finite quotient upstairs killing $\iota$.
The involution remains nonidentity after completion, since its
action on $H_1(\Sigma_2,\mathbb F_3)$ is $-I$. Thus the completed
kernel is exactly the finite group in (B.P1).

Put $P=\widehat\Gamma_{0,6}\triangleleft G$.
Its center is trivial: the full profinite Birman exact sequence,
with its free profinite point-pushing kernel, proves this by
induction from $\widehat\Gamma_{0,4}\cong\widehat F_2$.
Moreover
\[
                         C_G(P)=1.                       \tag{B.P3}
\]
Indeed an element centralizing $P$ has trivial permutation in
$G/P=S_6$, because the geometric $S_6$ factor injects into
$\operatorname{Out}(P)$ by
\cite[Corollary 2.8]{CA_HMM}. It then lies in $P$ and its center.
Consequently $Z(G)=1$, and (B.P1) gives
$Z(\widehat\Gamma_2)=\langle\iota\rangle$.
The kernel in (B.P1) is therefore characteristic, so $\alpha_F$
induces an automorphism $\bar\alpha_F$ of $G$.

The subgroup $P$ is the closed normal subgroup generated by the
$\sigma_i^2$. To see the exact quotient, add the relations
$\sigma_i^2=1$ to the sphere braid presentation: the braid and
commutation relations become the Coxeter presentation of $S_6$,
and both sphere relations already hold there. The action A.6
sends each $T_i$ to a conjugate of $T_i^{\lambda(F)}$.
It follows that $\bar\alpha_F$ sends every $\sigma_i^2$ into $P$;
the same argument for $F^{-1}$ proves equality of the image with
$P$. This proves pure-subgroup preservation through the actual
normal-generation map, without a separate characteristic-level
assertion.

Choose the separating circle which cuts the surface into two
one-holed tori and whose quotient separates branch labels
$\{1,2,3\}$ from $\{4,5,6\}$. Choose a quilt compatible with this
cut. The sewing compatibility proved in A.6 makes
\[
 K_2=\langle T_1,T_2\rangle\cong\widehat B_3
\]
invariant; its restriction is the one-holed-torus action with
parameter $F$, up to an inner automorphism of $K_2$.
More explicitly, for $a=T_1$, $b=T_2$, $F=(\lambda,f)$, the local
formula preceding (TA.8) is
\[
\begin{aligned}
 \Phi_F(a)&=a^\lambda,\\
 \Phi_F(b)&=a^{r_F}f(b^2,a^2)b^\lambda
                         f(b^2,a^2)^{-1}a^{-r_F},\\
 r_F&=N(\lambda-1)/2-8c_2(F),
\end{aligned}
\]
where $N$ is the fixed quilt adjustment integer. Antisymmetry gives
$f(b^2,a^2)=f(a^2,b^2)^{-1}$, so this is exactly the standard
three-strand action with parameter $F$, conjugated by $a^{r_F}$.
Thus its parameter and its entire inner correction are specified.
Downstairs,
$K=q(K_2)=\langle\sigma_1,\sigma_2\rangle$ is the braid subgroup
supported in the disk containing the first three branch points.
The actual profinite embedding and the commuting
one-holed-torus/hyperelliptic diagram are also given in
\cite[Lemma 3.2 and its proof]{ET_NS2000}.
For completeness, its injectivity can be checked through the maps
below: forgetting two complementary branch points sends it to
$\widehat B_3/Z(\widehat B_3)$; its remaining central full twist is
the twist about the essential three--three separating circle and
has an injected procyclic inertia group. Thus no kernel is lost
on full completion.

Hoshi--Minamide--Mochizuki give a unique pair
$(\Theta,\sigma)\in\widehat{\mathrm{GT}}\times S_6$ for the
outer restriction of $\bar\alpha_F$ to $P$.
Here is a finite pure-label test proving that $\sigma$ preserves
the first triple. Use the literal quotient
\[
\begin{aligned}
 V=H_1(P,\mathbb F_2)
  &=\left(\bigoplus_{1\leq i<j\leq5}\mathbb F_2 e_{ij}\right)
       \Big/\left\langle\sum_{i<j\leq5}e_{ij}\right\rangle,\\
 e_{i6}&=\sum_{j\leq5,\ j\ne i}e_{ij}.
\end{aligned}                                             \tag{B.P4}
\]
This is the abelianization of $P_5/Z(P_5)$ modulo two; $e_{ij}$ is
the class of the positive twist enclosing precisely branch points
$i,j$. Put $W=\langle e_{12},e_{13},e_{23}\rangle$.
It has dimension three. In $V/W$ the six classes
$e_{14},e_{15},e_{24},e_{25},e_{34},e_{35}$ form a basis, and
\[
\begin{aligned}
 e_{45}&=e_{14}+e_{15}+e_{24}+e_{25}+e_{34}+e_{35},\\
 e_{16}&=e_{14}+e_{15},&e_{26}&=e_{24}+e_{25},
       &e_{36}&=e_{34}+e_{35},\\
 e_{46}&=e_{15}+e_{25}+e_{35},
       &e_{56}&=e_{14}+e_{24}+e_{34}.
\end{aligned}                                                   \tag{B.P5}
\]
These fifteen pair classes are distinct and nonzero; among them
exactly $e_{12},e_{13},e_{23}$ belong to $W$.
The two adjacent squares in $K$ are sent by its
one-holed-torus action to unit powers of pair twists conjugated
inside $K$, so their images in $V$ belong to $W$.
The GT factor $\Theta$ acts on these classes by its odd multiplier,
hence trivially over $\mathbb F_2$; the other factor sends
$e_{ij}$ to $e_{\sigma(i)\sigma(j)}$.
It follows that both $\sigma\{1,2\}$ and $\sigma\{2,3\}$ lie
in the first triple, and hence
$\sigma\{1,2,3\}=\{1,2,3\}$.
The exact program \path{check_hyperelliptic_label_test.py} checks
(B.P4)--(B.P5) and all $720$ permutations; the allowed permutations
are exactly the $36$ elements of $S_3\times S_3$.

Realize $\sigma$ by a discrete braid $h$ supported in the first
triple together with a discrete braid supported in the complementary
triple. The latter commutes with $K$; the former acts on $K$ by
an inner automorphism. Thus
$\alpha_1=\operatorname{Inn}(h^{-1})\bar\alpha_F$ still preserves
$K$ and still has outer parameter $F$ there. Its outer action on
$P$ is exactly $\Theta$. For the standard ordinary genus-zero
representative $\gamma_\Theta$, this means that there is a single
$u\in P$ with
$\alpha_1|_P=\operatorname{Inn}(u)\gamma_\Theta|_P$.
This equality extends to every element of $G$: two automorphisms
agreeing on $P$ differ on any $g\in G$ by an element centralizing
$P$, and that centralizer is trivial by (B.P3). Therefore
\[
                  \alpha_1=\operatorname{Inn}(u)\gamma_\Theta
                         \quad\hbox{on }G.                \tag{B.P6}
\]

Let $G_0\subset G$ be the open mixed subgroup which permutes
$1,2,3$ and fixes $4,5,6$ individually. It contains both $K$ and
$u$. The actual forgetting map retaining $1,2,3,6$ is
\[
 \pi:G_0\longrightarrow
       \widehat\Gamma_{0,[3],1}
       \cong\widehat{\operatorname{PSL}_2(\mathbb Z)}
       \cong\widehat{C_2*C_3}.                            \tag{B.P7}
\]
Its restriction to $K$ sends the two positive half-twists to the
positive generators $a,b$ and kills exactly the central full twist
$(ab)^3$. This is the same positive modular marking as the
one-holed-torus cap followed by its central involution quotient:
\[
 a\longmapsto
   \left[\begin{smallmatrix}1&1\\0&1\end{smallmatrix}\right],
 \qquad
 b\longmapsto
   \left[\begin{smallmatrix}1&0\\-1&1\end{smallmatrix}\right],
 \qquad s=aba,\quad r=ab,\quad a=(sr)^{-1}.
\]
The equality of the two routes is literal on these generating
twists. The genus-zero braid formula restricts to the first three
strands, so $\pi\gamma_\Theta|_K$ has parameter $\Theta$ in this
marking. Applying $\pi$ to (B.P6) on $K$ shows that its outer action
equals the one with parameter $F$; the inner witness is the actual
element $\pi(u)$, together with the already retained inner change
inside $K$.

The multipliers of $F$ and $\Theta$ agree as well. For example,
the image of $\sigma_1^2$ has integral abelianization
$\lambda(F)e_{ij}$ for the pair determined by its actual conjugator.
The HMM description gives
$\lambda(\Theta)e_{\sigma(1)\sigma(2)}$. Reducing this equality
modulo two and using the fifteen distinct pair vectors just proved
identifies the two pairs. Their common vector is primitive in the
integral abelianization of $P_5/Z(P_5)$, so the equality over
$\widehat{\mathbb Z}$ forces
$\lambda(F)=\lambda(\Theta)$. Their quotient therefore has multiplier one.
The faithful normalized multiplier-one map to
$\operatorname{Out}(\widehat{C_2*C_3})$ proved after (B.M13) now
gives $\Theta=F$. Explicitly, an inner equality of two normalized
representatives first has a pure inner witness, by the faithful
$S_3$ action on the pure abelianization. Fixing the primitive cusp
$x$ makes that witness $x^t$. The second generator then gives
$f_2=y^c f_1x^{-t}$; zero abelianization of both parameters forces
$c=t=0$. This supplies the exact parameter-detection morphism.

We have proved $\bar\alpha_F=\operatorname{Inn}(hu)\gamma_F$.
Surjectivity of (B.P1) gives an actual lift $v$ of $hu$ and hence
the literal square (B.P2). Other quilts change the initial action
by the explicit inner comparisons of A.6 and so give the same
outer assertion.
\end{proof}
